\mxunit{premiere-c01}{Suites numériques et modèles discrets}\label{C01}
\begin{keypoint}[Le parcours du chapitre]
\textbf{Objectifs.} Passer du récit à une suite ; distinguer accroissement et taux ; calculer des termes et des sommes ; justifier une variation ; expérimenter un seuil.

\textbf{Prérequis.} Calcul littéral, puissances entières, pourcentages, fonction affine, boucle Python.

\textbf{Repère annuel.} 8–17 septembre : 8 h, puis reprises\end{keypoint}
\section*{Diagnostic : six questions, 10 minutes}\label{C01-D}
Répondre sans consulter le cours. Indiquer les questions pour lesquelles une aide est nécessaire. Ce diagnostic ne donne pas lieu à une note.
\begin{enumerate}
\item Calculer $3+2\times4$.
\item Calculer $2^0$, $2^3$ et $2^{-2}$.
\item Un effectif de 80 augmente de $5\,\%$. Donner le nouvel effectif.
\item Pour $f(x)=3x-2$, calculer $f(0)$ et $f(4)$.
\item Combien d’entiers y a-t-il de 3 à 10, bornes comprises ?
\item Une boucle \code{for k in range(4)} prend quelles valeurs de $k$ ?
\end{enumerate}
\paragraph{Après le diagnostic.} Reprendre les notions signalées avec le professeur ; refaire une question analogue, puis expliquer la règle utilisée.
\paragraph{Oral et logique.} En 90 secondes, expliquer la différence entre rang, terme et nombre de termes. Donner un contre-exemple à une règle devinée sur trois valeurs.
\mxnewpage
\section*{Chercher avant de formaliser}\label{C01-RECH}
En binômes, comparer deux offres de collecte : 20 kg la première semaine puis 5 kg de plus chaque semaine ; 20 kg puis une hausse de 20 \% chaque semaine. Conjecturer laquelle collecte le plus au total en six semaines, puis justifier. Prévoir 10 min de recherche et 5 min de confrontation.
\paragraph{Organisation.} Recherche individuelle 5 min, confrontation à deux 5 min, mise en commun 5 min. Garder une trace d’un essai, même non abouti. Une conjecture doit ensuite être justifiée.
\section{Définir et représenter une suite}\label{C01-U1}
\begin{keypoint}[Définition et résultats]
Une suite réelle est une fonction dont la variable est un entier : à chaque rang $n$ autorisé correspond un réel $u_n$. La notation $(u_n)$ désigne la suite, $u_n$ un terme. Une formule explicite donne directement $u_n$ en fonction de $n$. Une définition par récurrence fournit un premier terme et une règle de passage. La représentation est un ensemble de points $(n;u_n)$ : on ne relie pas nécessairement ces points.
\end{keypoint}
\paragraph{Statut et justification.} Définition. L’unicité des termes obtenus à partir d’un premier terme et d’une règle déterministe est admise ; aucun raisonnement par récurrence formel n’est requis.
\begin{example}Pour $u_n=n^2+1$, $u_0=1$, $u_3=10$. Pour $v_0=1$ et $v_{n+1}=2v_n+1$, $v_1=3$, $v_2=7$, $v_3=15$. On ne peut pas remplacer $v_3$ par $2\times3+1$ : la règle porte sur le terme précédent.\end{example}\label{C01-EX1}
\mxinclude{figures/C01-U1}
\begin{pitfall}Le rang n’est pas la valeur. Une liste finie de termes ne détermine pas une règle unique.\end{pitfall}
\begin{application}Calculer les quatre premiers termes de $w_0=2$, $w_{n+1}=w_n+2n+1$ ; proposer une formule explicite.\end{application}\label{C01-AP1}
\section{Suites arithmétiques}\label{C01-U2}
\begin{keypoint}[Définition et résultats]
Une suite est arithmétique de raison $r$ si $u_{n+1}=u_n+r$ pour tout rang autorisé. Alors $u_n=u_p+(n-p)r$ pour $n\geq p$. Si $r>0$, elle est strictement croissante ; si $r<0$, strictement décroissante ; si $r=0$, constante. Les points sont alignés sur une droite.
\end{keypoint}
\paragraph{Statut et justification.} Pour $p$ fixé, la suite $v_n=u_p+(n-p)r$ vérifie $v_p=u_p$ et $v_{n+1}=v_n+r$. Par unicité des termes successifs, $u_n=v_n$. La variation résulte du signe de $u_{n+1}-u_n=r$.
\begin{example}Une réserve contient 140 L et reçoit 12 L par minute. Après $n$ minutes, $u_n=140+12n$. À 5 min : 200 L. Si l’on commence l’indexation à 1 avec 140 L, on écrit $u_n=140+12(n-1)$.\end{example}\label{C01-EX2}
\begin{pitfall}Ajouter un pourcentage constant ne donne généralement pas une suite arithmétique.\end{pitfall}
\begin{application}On sait $u_3=17$ et $u_8=37$ pour une suite arithmétique. Déterminer $r$ et $u_0$.\end{application}\label{C01-AP2}
\section{Suites géométriques}\label{C01-U3}
\begin{keypoint}[Définition et résultats]
Une suite est géométrique de raison $q$ si $u_{n+1}=qu_n$. Pour $n\geq p$, $u_n=u_pq^{n-p}$, avec la convention qu’un produit vide vaut 1. Un taux constant $t$ donne $q=1+t$. La variation se décide avec $u_{n+1}-u_n=(q-1)u_n$. Pour $q<0$, les signes peuvent alterner.
\end{keypoint}
\paragraph{Statut et justification.} La formule donne le premier terme et sa multiplication par $q$ donne le terme suivant. L’unicité de la génération justifie la formule. Pour $q=0$, tous les termes après le premier sont nuls ; traiter ce cas directement.
\begin{example}Une quantité de 160 mg perd 25 \% à chaque étape : $u_n=160(\dfrac{3}{4})^n$. Les termes sont positifs et $u_{n+1}-u_n=-\dfrac{u_n}{4}<0$. On obtient $u_2=90$ mg, sans confondre baisse de 25 \% et retrait de 25 mg.\end{example}\label{C01-EX3}
\begin{pitfall}$q>1$ ne suffit pas à conclure que toute suite géométrique est croissante : le signe du premier terme compte.\end{pitfall}
\begin{application}Étudier la variation de $u_n=-2\times3^n$, puis de $v_n=2(-1)^n$.\end{application}\label{C01-AP3}
\section{Sommes, indices et démonstrations}\label{C01-U4}
\begin{keypoint}[Définition et résultats]
Pour $n\geq1$, $1+2+\cdots +n=\dfrac{n(n+1)}{2}$. Pour $q\ne1$ et $n\geq0$, $1+q+\cdots +q^n=\dfrac{1-q^{n+1}}{1-q}$. Pour $q=1$, la somme vaut $n+1$. Ainsi, de $u_p$ à $u_n$, une somme arithmétique vaut $\dfrac{(n-p+1)(u_p+u_n)}{2}$ ; une somme géométrique vaut $\dfrac{u_p(1-q^{n-p+1})}{1-q}$ si $q\ne1$.
\end{keypoint}
\paragraph{Statut et justification.} En écrivant $S=1+\cdots+n$ dans les deux sens, chaque paire de la somme $2S$ vaut $n+1$, avec $n$ paires. Pour $T=1+q+\cdots+q^n$, soustraire $qT=q+\cdots+q^{n+1}$ : $(1-q)T=1-q^{n+1}$. Le cas $q=1$ doit précéder la division.
\begin{example}Pour $u_n=3+2n$, $u_2+\cdots +u_6=\dfrac{5(7+15)}{2}=55$. Pour $v_n=2\times3^n$, $v_0+\cdots +v_3=\dfrac{2(1-3^4)}{1-3}=80$.\end{example}\label{C01-EX4}
\begin{pitfall}Compter les termes avant d’utiliser la formule : de 0 à n, il y en a n+1.\end{pitfall}
\begin{application}Calculer $\sum_{k=2}^{5}2^k$ sans symbole de sommation dans la réponse.\end{application}\label{C01-AP4}
\section{Conjecturer une limite et programmer un seuil}\label{C01-U5}
\begin{keypoint}[Définition et résultats]
Une limite est ici abordée intuitivement. Les termes $\dfrac{1}{n+1}$ se rapprochent de 0 ; $n^2$ grandit sans borne ; $(-1)^n$ ne se stabilise pas. Une expérience numérique suggère une conclusion mais ne prouve pas un comportement pour tous les rangs. Un algorithme de seuil doit préciser le sens de l’inégalité et prévoir un arrêt.
\end{keypoint}
\paragraph{Statut et justification.} Constats illustratifs, sans théorie formelle des limites. La correction d’un seuil entier se vérifie en contrôlant le rang trouvé et le rang précédent, lorsque la monotonie est établie.
\begin{example}Pour $u_n=80(\dfrac{1}{2})^n$, le premier rang tel que $u_n<6$ est 4 : $u_3=10\geq6$ et $u_4=5<6$. Le programme commence avec $n=0,u=80$ et répète la division par 2 tant que $u\geq6$.\end{example}\label{C01-EX5}
\begin{pitfall}Une boucle sans borne peut ne jamais s’arrêter si le seuil est inaccessible ; $<$ et $\leq$ ne sont pas interchangeables.\end{pitfall}
\begin{application}Pour $u_n=48(\dfrac{1}{2})^n$, déterminer les premiers rangs où $u_n\leq6$ et où $u_n<6$.\end{application}\label{C01-AP5}
\section*{Parcours d’exercices}\label{C01-PARC}
Commencer par les exercices 1 et 2. Choisir ensuite un exercice 3 ou 4 sans aide. Les exercices 5 à 7 sont une réserve de transfert, de recherche ou de programmation : le professeur en désigne un. L’exercice 8 se reprend à distance.
\begingroup
\mxexercisehead{C01-E1}{Prendre en main les indices}
On considère la suite $(u_n)$ définie pour tout entier naturel $n$ par $u_n=2n^2-3n$. Le premier terme est donc celui de rang 0.
\begin{enumerate}[label=\arabic*.]
\item Calculer séparément $u_0$, $u_1$ et $u_4$, en faisant apparaître la substitution de $n$.
\item Exprimer $u_{n+1}$ en fonction de $n$, puis développer et réduire cette expression.
\item Calculer $u_{n+1}-u_n$. En étudiant son signe pour les rangs entiers, décrire les variations de la suite.
\end{enumerate}
\mxexercisehead{C01-E2}{Choisir le bon modèle}
Une pépinière possède aujourd’hui 100 jeunes plants. Elle envisage deux modèles pour prévoir son effectif : dans le modèle A, l’effectif augmente de 15 plants par mois ; dans le modèle B, il augmente de 10 pour cent par mois. On note $a_n$ et $b_n$ les effectifs prévus après $n$ mois, avec $a_0=b_0=100$. Les deux modèles sont étudiés séparément.
\begin{enumerate}[label=\arabic*.]
\item \textbf{Modèle A.} Écrire la relation entre $a_{n+1}$ et $a_n$. Reconnaître la nature de la suite, puis exprimer $a_n$ en fonction de $n$. Calculer $a_2$.
\item \textbf{Modèle B, après l’étude des suites géométriques.} Reprendre les mêmes questions pour $(b_n)$ et calculer $b_2$.
\item Pour chaque modèle, expliquer ce que l’on suppose constant. Citer une raison pour laquelle une prévision à très long terme pourrait devenir irréaliste.
\end{enumerate}
\mxexercisehead{C01-E3}{Retrouver une suite}
La suite $(u_n)$ est géométrique, définie à partir du rang 0. Sa raison $q$ est strictement positive. On connaît seulement $u_1=6$ et $u_3=54$.
\begin{enumerate}[label=\arabic*.]
\item Déterminer $q$ en reliant les deux termes connus. Présenter toutes les étapes qui permettent de choisir sa valeur.
\item En déduire $u_0$, puis $u_5$.
\item Que changerait la suppression de l’hypothèse $q>0$ ? Donner toutes les raisons alors possibles et expliquer si la valeur de $u_5$ serait modifiée.
\end{enumerate}
\mxexercisehead{C01-E4}{Des rangées de sièges}
Une salle de spectacle comporte 12 rangées. La première rangée comprend 18 sièges ; chaque rangée suivante comprend deux sièges de plus que la précédente. On note $u_n$ le nombre de sièges de la rangée $n$, pour $1\leq n\leq12$.
\begin{enumerate}[label=\arabic*.]
\item Définir la suite $(u_n)$ et calculer le nombre de sièges de la dernière rangée.
\item Calculer la capacité totale de la salle. Expliquer le nombre de termes retenu dans la somme.
\item Calculer le nombre moyen de sièges par rangée et utiliser ce résultat pour contrôler le total trouvé.
\end{enumerate}
\mxexercisehead{C01-E5}{Accumuler des carrés}
On construit des motifs avec des carreaux carrés identiques. Le motif de rang 1 contient un carreau. Pour passer du motif de rang $n$ à celui de rang $n+1$, on ajoute $2n+1$ carreaux. On note $u_n$ le nombre de carreaux du motif de rang $n$, pour $n\geq1$.
\begin{enumerate}[label=\arabic*.]
\item Calculer $u_2$ et $u_3$. Dessiner une disposition de ces carreaux qui suggère une formule pour $u_n$.
\item Conjecturer une expression de $u_n$ en fonction de $n$.
\item Vérifier que cette expression donne le premier terme et satisfait la règle de construction. Expliquer pourquoi ces deux vérifications identifient la suite.
\end{enumerate}
\mxinclude{figures/C01-E5}
\mxexercisehead{C01-E6}{Deux remboursements}
Une personne doit rembourser 600 unités monétaires, sans intérêts. Deux plans sont envisagés. Avec le plan A, elle verse 50 unités à la fin de chaque mois jusqu’au remboursement. Avec le plan B, elle verse chaque mois la moitié de ce qu’elle doit encore. Les calculs du modèle B sont effectués sans arrondi. On note $a_n$ et $b_n$ les dettes restantes après le versement du mois $n$, avec $a_0=b_0=600$.
\begin{enumerate}[label=\arabic*.]
\item Exprimer $a_n$ jusqu’à l’extinction de la dette. À quel mois le remboursement est-il achevé ?
\item Exprimer $b_n$. Déterminer le premier mois où la dette restante est strictement inférieure à 10 unités ; vérifier le mois précédent.
\item Dans le modèle B sans arrondi, la dette devient-elle exactement nulle après un nombre fini de mois ? Justifier et commenter la différence avec une situation réelle.
\end{enumerate}
\mxexercisehead{C01-E7}{Lire une liste}
Un programme doit conserver les termes de la suite définie par $u_0=2$ et $u_{n+1}=3u_n-1$. Il commence par les instructions \code{u=2} et \code{L=[u]}. Il exécute ensuite cinq fois, dans cet ordre, \code{u=3*u-1} puis \code{L.append(u)}.
\begin{enumerate}[label=\arabic*.]
\item Exécuter les deux premiers passages de la boucle à la main. Donner alors le contenu de la liste.
\item Combien d’éléments contient la liste à la fin des cinq passages ? Quel est le rang du dernier terme enregistré ?
\item Écrire le programme complet en Python, puis vérifier la valeur du dernier élément. Expliquer pourquoi le nombre de termes dépasse d’une unité le nombre de passages.
\end{enumerate}
\mxexercisehead{C01-E8}{Reprise différée}
Cette question se traite une semaine après le chapitre, sans consulter les formules. Une suite est définie pour $n\geq2$ par $u_2=12$ et $u_{n+1}=u_n-3$.
\begin{enumerate}[label=\arabic*.]
\item Calculer $u_7$ en précisant combien de passages séparent les rangs 2 et 7.
\item Calculer la somme $u_2+u_3+\cdots+u_7$. Expliquer la différence entre le nombre de termes de cette somme et le nombre de passages calculé à la première question.
\end{enumerate}
\paragraph{Travail personnel.} Effectuer uniquement la tâche désignée à l’issue de la séance, dans la progression datée : 10–20 min maximum. Les exercices dépendent des notions déjà étudiées ; certaines questions sont donc différées. Reprendre l’exercice 8 une semaine après le chapitre (10 min). Si le temps prévu est dépassé, noter la difficulté et apporter les essais ; ne pas accumuler toute la réserve.
\endgroup
\mxneed{12}\section*{À faire dans le manuel CQFD}
Ces exercices se traitent dans votre manuel. La progression indique la séance ; ils remplacent le travail personnel prévu ce jour-là (10 à 20 minutes au total).
\paragraph{CQFD, exercice 9 p. 64, questions a et b.}\label{C01-M1} Calculer les quatre premiers termes. Écrire les substitutions avant de simplifier.
\paragraph{CQFD, exercice 11 p. 64, questions toutes.}\label{C01-M2} Calculer les termes successivement et conserver la trace des quatre étapes.
\begingroup\small
\mxnewpage\section*{QCM 1 : comprendre et choisir}\label{C01-Q1}
Cinq questions pédagogiques à réponse unique. Choisir une réponse par question ; expliquer ensuite un distracteur. Ce format de chapitre ne reproduit pas le format officiel de l’EAM.
\begin{enumerate}
\item $u_n=3n+2$ : $u_4$ vaut
\begin{enumerate}[label=\Alph*.,itemsep=2pt]
\item $14$
\item $18$
\item $12$
\item $6$
\end{enumerate}
\item $u_0=2$, $u_{n+1}=u_n+5$ : $u_3$ vaut
\begin{enumerate}[label=\Alph*.,itemsep=2pt]
\item $10$
\item $12$
\item $17$
\item $15$
\end{enumerate}
\item Une hausse de 8 \% donne la raison
\begin{enumerate}[label=\Alph*.,itemsep=2pt]
\item $0{,}08$
\item $1{,}8$
\item $0{,}92$
\item $1{,}08$
\end{enumerate}
\item De $u_3$ à $u_9$ inclus, on compte
\begin{enumerate}[label=\Alph*.,itemsep=2pt]
\item $6$
\item $7$
\item $9$
\item $12$
\end{enumerate}
\item $1+2+\cdots+20$ vaut
\begin{enumerate}[label=\Alph*.,itemsep=2pt]
\item $200$
\item $420$
\item $210$
\item $190$
\end{enumerate}
\end{enumerate}
\endgroup
\begingroup\small
\mxneed{18}\section*{QCM 2 : comprendre et choisir}\label{C01-Q2}
Cinq questions pédagogiques à réponse unique. Choisir une réponse par question ; expliquer ensuite un distracteur. Ce format de chapitre ne reproduit pas le format officiel de l’EAM.
\begin{enumerate}
\item $u_n=-3\times2^n$ est
\begin{enumerate}[label=\Alph*.,itemsep=2pt]
\item strictement croissante
\item constante
\item ni croissante ni décroissante
\item strictement décroissante
\end{enumerate}
\item $1+2+4+8$ vaut
\begin{enumerate}[label=\Alph*.,itemsep=2pt]
\item $15$
\item $16$
\item $14$
\item $8$
\end{enumerate}
\item Pour $u_n=\dfrac{16}{2^n}$, premier rang où $u_n<2$ :
\begin{enumerate}[label=\Alph*.,itemsep=2pt]
\item $2$
\item $3$
\item $4$
\item $5$
\end{enumerate}
\item La suite $(-1)^n$
\begin{enumerate}[label=\Alph*.,itemsep=2pt]
\item tend vers 0
\item ne possède pas de limite
\item est croissante
\item est constante
\end{enumerate}
\item Une liste contient $u_0$ à $u_{10}$ ; sa longueur est
\begin{enumerate}[label=\Alph*.,itemsep=2pt]
\item $10$
\item $9$
\item $12$
\item $11$
\end{enumerate}
\end{enumerate}
\subsection*{Vrai ou faux ? Justifier.}\label{C01-VF}
\begin{enumerate}
\item Toute suite géométrique de raison $q>1$ est croissante.
\item Si deux suites ont leurs trois premiers termes égaux, elles sont égales.
\item Pour $q=1$, $1+q+\cdots+q^n=n+1$.
\end{enumerate}
\begin{keypoint}[Bilan personnel]Je distingue explicite et récurrente ; je reconnais un accroissement ou un taux constant ; je compte les termes ; je contrôle le rang précédent d’un seuil.

Pour chaque objectif, indiquer : acquis sans aide ; réussi avec aide ; à retravailler. Écrire un exemple et une erreur à éviter.\end{keypoint}
\endgroup
\ifmxteacher\mxinclude{corriges/C01}\fi
